\section{Implementación de Filtros Deterministas}

En los sistemas contemporáneos de composición automatizada e inteligencia artificial aplicada a la música, uno de los desafíos teóricos y prácticos más significativos radica en la conciliación del comportamiento estocástico y creativo de los modelos de aprendizaje profundo con las restricciones normativas y formales descritas por la teoría musical clásica. Las redes neuronales profundas, tales como las redes de memoria a corto y largo plazo (LSTM) o los modelos basados en mecanismos de atención (Transformers), demuestran una capacidad sobresaliente para aprender representaciones complejas, transiciones de probabilidad y estilos a partir de corpus masivos de datos. Sin embargo, al estar fundamentadas en optimizaciones probabilísticas y de mínimos cuadrados (tales como la pérdida del error cuadrático medio, MSE), estas redes carecen intrínsecamente de un mecanismo determinista para garantizar el cumplimiento de restricciones rígidas (`hard constraints`). 

Por ejemplo, una red neuronal entrenada para la generación melódica puede sugerir un intervalo de novena menor o una nota extremadamente disonante fuera de la escala base simplemente porque la distribución de probabilidad local del modelo de lenguaje musical presentaba un residuo estocástico propicio en dicho punto. Si bien este comportamiento puede ser tolerado en contextos de improvisación libre o música contemporánea no tonal, resulta inaceptable bajo los marcos de la armonía tonal académica y el contrapunto estricto, donde errores como la resolución incorrecta de disonancias o la existencia de saltos melódicos no cantables invalidan la estructura de la pieza.

Para resolver esta limitación, el presente trabajo introduce una arquitectura híbrida de carácter neurosimbólico, estructurada en múltiples capas. En esta arquitectura, el componente de aprendizaje profundo (la red neuronal generativa parametrizada) actúa como un motor creativo de exploración libre que propone secuencias de notas candidatas con base en el estilo aprendido del conjunto de datos. Simultáneamente, el componente simbólico, constituido por un filtro determinista estructurado como un árbol de decisión parametrizable, interviene en la salida del modelo generativo. Este árbol actúa como un validador normativo de tipo "`hard constraint satisfaction`" (satisfacción de restricciones duras) que evalúa y modifica localmente la secuencia musical antes de su codificación en formato MIDI.

El ciclo dinámico de inferencia y post-procesamiento híbrido se formaliza de la siguiente manera. Sea una melodía candidata generada por la red neuronal expresada como una matriz $\mathbf{M} \in \mathbb{R}^{N \times 4}$, donde cada fila $j \in \{0, 1, \dots, N-1\}$ representa un evento musical compuesto por el vector:
\begin{equation}
    \mathbf{v}_j = [pc_j, d_j, deg_j, int_j]
\end{equation}
Donde $pc_j$ es la clase de altura (pitch class), $d_j$ es la duración temporal en unidades relativas de semicorchea, $deg_j$ representa el grado relativo a la escala activa, y $int_j$ es el intervalo en tonos respecto a la nota inmediatamente anterior. La matriz pasa a través de la función de filtrado determinista del árbol de decisión $f_{\text{árbol}}: \mathbb{R}^{N \times 4} \to \mathbb{R}^{N \times 4}$, la cual evalúa y corrige las desviaciones respecto a los límites teóricos parametrizados por el usuario. 

A diferencia de los enfoques tradicionales de generación restringida, los cuales descartan por completo las secuencias que no cumplen con los requisitos mínimos (lo que genera un costo computacional inasumible debido a la continua re-generación estocástica), el árbol de decisión propuesto opera bajo una estrategia de "reparación local". El árbol identifica el punto exacto de la desviación normativa y aplica una corrección matemática directa sobre la nota conflictiva y sus vecindades. Esto preserva la coherencia temática del modelo neuronal mientras asegura el cumplimiento inequívoco de las leyes de la teoría armónica.

\subsection{Codificación de tonalidades, modos e inducción de mutaciones}

Para parametrizar matemáticamente el comportamiento del filtro y del algoritmo genético, se requiere una formalización algebraica estricta del espacio de alturas y modos musicales. La clase de alturas (pitch class) se modela sobre el anillo cociente $\mathbb{Z}_{12} = \{0, 1, 2, \dots, 11\}$, donde cada entero representa una nota musical del temperamento igual, tomando a Do ($C$) como el origen de coordenadas $0$:
\begin{equation}
    C \equiv 0, \quad C\# \equiv 1, \quad D \equiv 2, \quad \dots, \quad B \equiv 11
\end{equation}

La tónica de la escala se codifica mediante un parámetro entero $T \in \{1, 2, \dots, 12\}$, el cual se transforma linealmente al espacio $\mathbb{Z}_{12}$ a través de la función de proyección:
\begin{equation}
    t_0 = (T - 1) \pmod{12}
\end{equation}
Esta tónica sirve como el origen tonal absoluto de la composición.

Los modos eclesiásticos o modos griegos diatónicos se codifican mediante un parámetro de modo $M \in \{1, 2, \dots, 7\}$. Cada modo $M$ se define matemáticamente como un conjunto ordenado de siete intervalos relativos $\mathcal{I}_M = \{i_0, i_1, \dots, i_6\} \subset \mathbb{Z}_{12}$ que representan las distancias en semitonos desde el centro tonal $t_0$. Los modos se estructuran algebraicamente según los siguientes vectores de intervalos:
\begin{align}
    \mathcal{I}_1 \text{ (Jónico)} &= \{0, 2, 4, 5, 7, 9, 11\} \\
    \mathcal{I}_2 \text{ (Dórico)} &= \{0, 2, 3, 5, 7, 9, 10\} \\
    \mathcal{I}_3 \text{ (Frigio)} &= \{0, 1, 3, 5, 7, 8, 10\} \\
    \mathcal{I}_4 \text{ (Lidio)} &= \{0, 2, 4, 6, 7, 9, 11\} \\
    \mathcal{I}_5 \text{ (Mixolidio)} &= \{0, 2, 4, 5, 7, 9, 10\} \\
    \mathcal{I}_6 \text{ (Eólico)} &= \{0, 2, 3, 5, 7, 8, 10\} \\
    \mathcal{I}_7 \text{ (Locrio)} &= \{0, 1, 3, 5, 6, 8, 10\}
\end{align}

Dado un par de parámetros $(T, M)$, la escala diatónica activa en el sistema se denota como el conjunto $\mathcal{D}_{T,M} \subset \mathbb{Z}_{12}$, definido formalmente como la traslación del conjunto de intervalos modales respecto a la tónica:
\begin{equation}
    \mathcal{D}_{T,M} = \{ (t_0 + i) \pmod{12} \mid i \in \mathcal{I}_M \}
\end{equation}

Una nota representada por una clase de altura $pc$ se define como diatónica si y solo si $pc \in \mathcal{D}_{T,M}$. En caso contrario, se define como una nota cromática o no diatónica.

En el marco evolutivo del sistema, el algoritmo genético interactivo opera sobre una población de cromosomas. Cada cromosoma es una matriz de melodía $\mathbf{M}$ de dimensiones $N \times 4$. La optimización de la red neuronal y la búsqueda de trayectorias melódicas estéticas dependen de los operadores genéticos de cruce y mutación aplicados a estas matrices. 

El operador de cruce (crossover) recibe dos matrices progenitoras $\mathbf{P}_A, \mathbf{P}_B \in \mathbb{R}^{N \times 4}$ y genera un hijo $\mathbf{H} \in \mathbb{R}^{N \times 4}$ mediante un corte de punto único. Sea $k \in \{1, 2, \dots, N-2\}$ un punto de corte aleatorio. El hijo resultante se construye mediante la concatenación vertical:
\begin{equation}
    \mathbf{H} = \begin{bmatrix} \mathbf{P}_A[0 \dots k-1] \\ \mathbf{P}_B[k \dots N-1] \end{bmatrix}
\end{equation}

El operador de mutación introduce variabilidad en el cromosoma con una probabilidad por gen de $p_m$. Para asegurar que las mutaciones generen individuos que posean sentido biológico y técnico para la red neuronal recurrente, las mutaciones se especializan de acuerdo a las características físicas del parámetro alterado:
\begin{enumerate}
    \item \textbf{Mutación de la Duración ($d_j$):} Si el parámetro rítmico de la nota en la posición $j$ es seleccionado para mutar, su valor se altera aleatoriamente sustituyéndose por un elemento del conjunto discreto de duraciones musicales válidas:
    \begin{equation}
        d_j \in \mathcal{U}_{\text{duraciones}} = \{0.25, 0.50, 0.75, 1.00, 1.50, 2.00, 3.00, 4.00\}
    \end{equation}
    Este conjunto discreto representa subdivisiones que van desde la semicorchea ($0.25$) hasta la redonda ($4.00$) en términos relativos, asegurando la consistencia rítmica de la partitura.
    
    \item \textbf{Mutación de la Altura ($pc_j$) con sentido estructural:} Si la altura de la nota en la posición $j$ muta, el sistema selecciona de manera uniforme un nuevo valor de clase de altura $pc'_j \in \mathcal{D}_{T,M}$. La asignación directa de una nota arbitraria destruiría la coherencia interválica local de la melodía y generaría saltos melódicos no deseados de gran amplitud. Para evitar este efecto, se calcula la diferencia en semitonos $\delta \in \mathbb{Z}$ entre la nueva nota y la nota antigua:
    \begin{equation}
        \delta = pc'_j - pc_j
    \end{equation}
    Esta diferencia se normaliza al intervalo diatónico más corto mediante:
    \begin{equation}
        \delta_n = \begin{cases} \delta - 12 & \text{si } \delta > 6 \\ \delta + 12 & \text{si } \delta \le -6 \\ \delta & \text{en otro caso} \end{cases}
    \end{equation}
    
    Para inducir la mutación de forma que "haga sentido" al codificador y al decodificador matricial, el algoritmo actualiza localmente tres variables de la matriz de forma simultánea:
    \begin{itemize}
        \item La clase de altura del gen mutado se actualiza: $pc_j \leftarrow pc'_j$.
        \item El grado de la escala se re-calcula usando la función de correspondencia diatónica: $deg_j \leftarrow g(pc'_j, t_0, \mathcal{I}_M)$.
        \item Se ajustan de forma balanceada las relaciones interválicas adyacentes para encapsular la mutación. El intervalo en tonos del evento actual se incrementa sumando la mitad de la desviación de semitonos:
        \begin{equation}
            int_j \leftarrow int_j + \frac{\delta_n}{2}
        \end{equation}
        Para evitar que esta alteración desplace de forma absoluta e indeseada todas las notas subsiguientes de la melodía (lo que arruinaría la progresión temática previa), se sustrae exactamente la misma cantidad al intervalo de la siguiente nota:
        \begin{equation}
            int_{j+1} \leftarrow int_{j+1} - \frac{\delta_n}{2}
        \end{equation}
        Si $j = N-1$ (la última nota de la secuencia), solo se altera $int_j$. Este mecanismo de conservación garantiza la preservación del contorno de alturas absoluto aguas abajo de la mutación.
    \end{itemize}
\end{enumerate}

\subsection{Lógica de poda y validación de secuencias candidatas}

El proceso de validación y poda de las melodías candidatas propuestas por el modelo estocástico se efectúa de manera secuencial a través de un árbol de decisión determinista. Este árbol analiza cada evento melódico y sus relaciones interválicas previas para detectar violaciones a las restricciones tonales y de conducción de voces.

En la teoría de conducción de voces y contrapunto, los intervalos musicales se clasifican rigurosamente por su amplitud física en semitonos absolutos $\Delta = |p_t - p_{t-1}|$:
\begin{itemize}
    \item \textbf{Paso o Grado Conjunto (Step):} Movimiento de $1$ o $2$ semitonos (segunda menor o segunda mayor). Representa la máxima suavidad melódica.
    \item \textbf{Salto Corto (Skip):} Movimiento de $3$ o $4$ semitonos (tercera menor o tercera mayor). Introduce una tensión moderada pero fluida.
    \item \textbf{Salto Largo (Leap):} Movimiento de $5$ o más semitonos (cuarta justa, quinta justa o mayores). Introduce una gran tensión y debe resolverse de forma contraria para mantener la cantabilidad.
\end{itemize}

El árbol de decisión opera sobre la secuencia melódica aplicando un algoritmo recursivo de validación. La lógica interna de este árbol se describe mediante la secuencia de pruebas lógicas y ramificaciones condicionales estructuradas a continuación:

\begin{algorithm}[H]
\caption{Algoritmo del Árbol de Decisión para Corrección de Melodías}
\begin{algorithmic}[1]
\REQUIRE Matriz de melodía original $\mathbf{M}$, parámetros $T, M, \text{max\_skips}, \text{max\_leaps}, \text{max\_non\_diatonic}$
\ENSURE Matriz de melodía filtrada y corregida $\mathbf{M}'$
\STATE Obtener la escala diatónica activa $\mathcal{D}_{T,M}$ y la tónica $t_0$.
\STATE Inicializar contadores: $c_{\text{crom}} \leftarrow 0$, $c_{\text{skips}} \leftarrow 0$, $c_{\text{leaps}} \leftarrow 0$.
\STATE Reconstruir las alturas MIDI absolutas $P = \{p_0, p_1, \dots, p_{N-1}\}$ a partir de la columna de intervalos de $\mathbf{M}$.
\FOR{$i = 0$ \TO $N-1$}
    \STATE $pc \leftarrow p_i \pmod{12}$
    \COMMENT{\textbf{Nodo de Decisión 1: Diatonicidad}}
    \IF{$pc \notin \mathcal{D}_{T,M}$}
        \COMMENT{\textbf{Nodo de Decisión 2: Límite Cromático}}
        \IF{$c_{\text{crom}} \ge \text{max\_non\_diatonic}$}
            \STATE Modificar $p_i \leftarrow \text{get\_nearest\_diatonic}(p_i, \mathcal{D}_{T,M})$
            \STATE Actualizar $pc \leftarrow p_i \pmod{12}$
        \ELSE
            \STATE $c_{\text{crom}} \leftarrow c_{\text{crom}} + 1$
        \ENDIF
    \ENDIF
    \IF{$i > 0$}
        \STATE $\text{intervalo} \leftarrow p_i - p_{i-1}$
        \STATE $\Delta \leftarrow |\text{intervalo}|$
        \COMMENT{\textbf{Nodo de Decisión 3: Clasificación de Saltos}}
        \IF{$\Delta \in \{3, 4\}$}
            \COMMENT{\textbf{Nodo de Decisión 4: Límite de Skips}}
            \IF{$c_{\text{skips}} \ge \text{max\_skips}$}
                \STATE $\text{dirección} \leftarrow \text{signo}(\text{intervalo})$
                \STATE $\text{step\_size} \leftarrow \text{diatonic\_step}(p_{i-1}, \text{dirección}, \mathcal{D}_{T,M})$
                \STATE $p_i \leftarrow p_{i-1} + \text{dirección} \times \text{step\_size}$
                \STATE Asegurar $p_i$ diatónico.
            \ELSE
                \STATE $c_{\text{skips}} \leftarrow c_{\text{skips}} + 1$
            \ENDIF
        \ELSIF{$\Delta \ge 5$}
            \COMMENT{\textbf{Nodo de Decisión 5: Límite de Leaps}}
            \IF{$c_{\text{leaps}} \ge \text{max\_leaps}$}
                \STATE $\text{dirección} \leftarrow \text{signo}(\text{intervalo})$
                \COMMENT{\textbf{Intento de degradación suave a Skip}}
                \IF{$c_{\text{skips}} < \text{max\_skips}$}
                    \STATE $p_i \leftarrow p_{i-1} + \text{dirección} \times 3$
                    \IF{$p_i \pmod{12} \notin \mathcal{D}_{T,M}$}
                        \STATE $p_i \leftarrow p_{i-1} + \text{dirección} \times 4$
                    \ENDIF
                    \STATE Asegurar $p_i$ diatónico.
                    \STATE $c_{\text{skips}} \leftarrow c_{\text{skips}} + 1$
                \ELSE
                    \COMMENT{\textbf{Forzar reducción estricta a Step}}
                    \STATE $\text{step\_size} \leftarrow \text{diatonic\_step}(p_{i-1}, \text{dirección}, \mathcal{D}_{T,M})$
                    \STATE $p_i \leftarrow p_{i-1} + \text{dirección} \times \text{step\_size}$
                    \STATE Asegurar $p_i$ diatónico.
                \ENDIF
            \ELSE
                \STATE $c_{\text{leaps}} \leftarrow c_{\text{leaps}} + 1$
            \ENDIF
        \ENDIF
    \ENDIF
\ENDFOR
\STATE Re-calcular los parámetros de la matriz $\mathbf{M}'$ basándose en las nuevas alturas corregidas $P$.
\RETURN $\mathbf{M}'$
\end{algorithmic}
\end{algorithm}

La función $\text{diatonic\_step}(pitch, dir, \text{escala})$ calcula la distancia al siguiente grado de la escala en la dirección indicada. Si la escala activa requiere un movimiento de $2$ semitonos, la función retorna $2$. Si el grado diatónico se encuentra a una distancia de segunda menor, retorna $1$.

Este mecanismo de poda asegura que si la red neuronal propone una melodía que salta abruptamente de forma repetida (ej. una sucesión de quintas y octavas disjuntas), el árbol de decisión interviene y "poda" la curva melódica, aproximando las notas a pasos diatónicos más cercanos. La validación se realiza de forma secuencial y cronológica para imitar el proceso humano de corrección y revisión melódica.

\subsection{Interfaz de parametrización: Tonalidad y modo}

El control creativo del sistema híbrido se canaliza mediante una interfaz de parametrización interactiva que actúa sobre los umbrales lógicos del árbol de decisión y los operadores evolutivos. Esta parametrización otorga al usuario un control exhaustivo sobre el nivel de disonancia, la agilidad de los saltos y el rigor estructural del material generado. 

La parametrización se gestiona de forma interactiva y permite regular las siguientes restricciones fundamentales:
\begin{table}[H]
    \centering
    \begin{tabular}{|l|c|l|}
        \hline
        \textbf{Parámetro} & \textbf{Rango de Entrada} & \textbf{Efecto en la Salida Melódica} \\ \hline
        $\text{tonic}$ & $1 \le T \le 12$ & Define el centro tonal absoluto de la composición. \\ \hline
        $\text{greek\_mode}$ & $1 \le M \le 7$ & Selecciona la distribución intervalar diatónica (modo griego). \\ \hline
        $\text{max\_skips}$ & $\ge 0$ & Controla la cantidad máxima de intervalos de tercera ($3$ o $4$ semitonos). \\ \hline
        $\text{max\_leaps}$ & $\ge 0$ & Controla la cantidad de intervalos disjuntos grandes ($\ge 5$ semitonos). \\ \hline
        $\text{max\_non\_diatonic}$ & $\ge 0$ & Regula la presencia de notas cromáticas o alteraciones accidentales. \\ \hline
    \end{tabular}
    \caption{Parámetros de la interfaz condicional del filtro armónico.}
\end{table}

La interacción de estos parámetros define el comportamiento del sistema en un continuo que va desde la rigidez académica hasta la experimentación libre:
\begin{itemize}
    \item \textbf{Configuración de Contrapunto Estricto:} Si el usuario parametriza $\text{max\_non\_diatonic} = 0$, $\text{max\_leaps} = 0$ y $\text{max\_skips} = 2$, el árbol de decisión podará de forma agresiva cualquier intento del modelo neuronal de realizar cromatismos o saltos melódicos pronunciados. El resultado será una melodía suave, principalmente por grados conjuntos, muy alineada con el estilo vocal del Renacimiento y las reglas de conducción clásicas.
    
    \item \textbf{Configuración de Expresividad Cromática Moderada:} Al relajar los parámetros a $\text{max\_non\_diatonic} = 3$, $\text{max\_skips} = 5$ y $\text{max\_leaps} = 2$, el árbol de decisión amplía su umbral de tolerancia. Esto permite la aparición controlada de notas accidentales de paso y saltos expresivos (ej. intervalos de quinta o sexta justa), dotando a la melodía de una estética más cercana al romanticismo o al jazz modal, sin llegar a colapsar en la atonalidad desorganizada.
\end{itemize}

Esta modularidad demuestra que la arquitectura neurosimbólica no limita la expresividad artística de la red neuronal. Por el contrario, la encauza mediante un conjunto dinámico de directrices técnicas controladas por el compositor. Esto permite al investigador explorar diferentes estilos musicales y regular de manera precisa el equilibrio entre la exploración probabilística del modelo recurrente y la explotación normativa del sistema de reglas simbólicas.
